% ============================================================
% Capítulo: Discusión (Bloque B.11)
% ============================================================
\chapter{Discusión}
\label{cap:discusion}

Este capítulo interpreta la evidencia del \autoref{cap:resultados} a la
luz de las cuatro preguntas de investigación planteadas en el
\autoref{cap:introduccion} -- sin repetir las cifras ya reportadas allí,
sino sintetizándolas y discutiendo lo que significan. Se cierra
comparando SGB-SaaS contra la literatura revisada en el
\autoref{cap:trabajos-relacionados}, señalando los hallazgos no buscados
que surgieron durante el desarrollo, y discutiendo las implicaciones
prácticas del proyecto para una biblioteca real.

\section{Respuesta a las preguntas de investigación}
\label{sec:disc-rqs}

\subsection*{RQ1 -- ¿En qué medida el caché reduce la latencia, y con qué margen se cumplen los umbrales?}

Ambos umbrales de rendimiento se cumplen con margen amplio -- p95
$19{,}49$\,ms contra un límite de $200$\,ms en caché caliente, p95
$7{,}50$\,ms contra $500$\,ms en frío
(\autoref{sec:res-rendimiento})--, así que en el sentido estricto que
exige la guía, la respuesta a RQ1 es afirmativa: el sistema cumple. Pero
la comparación \emph{entre} escenarios no permite concluir de forma
limpia que sea el caché específicamente el responsable de la diferencia
observada, porque el escenario \texttt{cache\_frio} tiene una limitación
metodológica real -- termina midiendo mayoritariamente consultas con
página vacía, no la consulta completa que \texttt{cache\_frio} debería
representar (\autoref{sec:res-rendimiento},
\autoref{cap:amenazas-validez}). Lo que sí se sostiene pese a esa
limitación es la dirección y el tamaño del efecto entre corridas de
\texttt{cache\_caliente}: Cliff's delta $=-0{,}68$ (grande, consistente en
las 5 corridas), exactamente la dirección que la hipótesis de caché
predeciría si la comparación fuera limpia. En síntesis: RQ1 se responde
con evidencia de cumplimiento de umbral sólida, y con evidencia de efecto
del caché sugestiva pero no aislada limpiamente -- una distinción que
este documento no difumina.

\subsection*{RQ2 -- ¿Qué proporción de controles OWASP presentó hallazgos reales, y cuántos se corrigieron?}

De los 6 controles auditados, 4 (A01, A05, A07, A09) presentaron un
hallazgo real -- el $66{,}7$\,\% reportado en
\autoref{sec:res-seguridad} -- y los 4 recibieron corrección verificable
dentro del mismo ciclo de desarrollo, re-confirmada después contra un
flujo de autenticación que cambió. Este es, en sí mismo, el dato más
relevante para RQ2: la mayoría de los controles auditados \textbf{sí}
tenían defectos reales en la implementación -- rol admin mal validado,
rate limiting de login ausente, logging de autenticación insuficiente, y
un código de estado HTTP incorrecto detectado durante la propia
verificación de A05 (\texttt{500} en vez de \texttt{404} en Swagger bajo
perfil \texttt{prod}). No fue una auditoría de casillas vacías que
confirma lo que ya se sabía: encontró y forzó a corregir defectos de
seguridad concretos. El control A03 (inyección) se verificó limpio, sin
hallazgo -- un resultado real también, no la ausencia de una prueba. A02
queda en un estado genuinamente mixto que RQ2 no puede resolver del todo
todavía: la decisión de arquitectura y la preparación del backend para
TLS están cerradas, pero la verificación final del header
\texttt{Strict-Transport-Security} depende de un despliegue público que
todavía no existe -- no es un hallazgo sin corregir, es una verificación
que no puede completarse sin una pieza de infraestructura externa al
propio código.

\subsection*{RQ3 -- ¿Cómo perciben usuarios reales la usabilidad del sistema (SUS)?}

RQ3 \textbf{no tiene respuesta todavía}. Con $N=0$ de los $N\geq15$
participantes exigidos, no existe ninguna puntuación SUS que interpretar
(\autoref{sec:res-usabilidad}). Forzar una respuesta aquí -- aunque fuera
calificada de ``preliminar'' o ``estimada'' -- sería exactamente el tipo
de afirmación no verificada que este documento se niega a producir en
cualquier otro capítulo; no hay razón para relajar ese criterio en la
discusión. La respuesta honesta a RQ3, a la fecha de este documento, es
que sigue siendo trabajo en progreso, bloqueado en la cadena de
dependencias correcta (protocolo de consentimiento ya definido
$\to$ despliegue público $\to$ reclutamiento $\to$ aplicación del
instrumento), no una omisión arbitraria.

\subsection*{RQ4 -- ¿Qué efecto tiene la estrategia híbrida CRUD-ORM/SP sobre trazabilidad y cobertura?}

La matriz de trazabilidad (\autoref{sec:matriz-resumen},
\autoref{cap:diseno-arquitectura}) muestra una distribución real y no
pareja: $48{,}8$\,\% de los requisitos son CRUD-ORM exclusivo,
$18{,}6$\,\% son procedimiento/función SQL exclusivo, y solo
$4{,}7$\,\% son genuinamente híbridos (requisitos transversales que
combinan ambos mecanismos) -- el resto no aplica un patrón de acceso a
datos directo. La cobertura de pruebas automatizadas, medida sobre esa
misma base de código, es del $81{,}64$\,\% de líneas
(\autoref{sec:res-cobertura}), muy por encima del umbral exigido, y no
se concentra de forma distinta entre el código que usa ORM y el que usa
SP -- ambos paquetes (\texttt{service}, que contiene la mayoría de las
invocaciones a SP, y el resto orientado a ORM) están sobre el $89$\,\%.
Con esta evidencia, RQ4 tiene una respuesta afirmativa en cuanto a
resultado: la estrategia híbrida no perjudicó ni la trazabilidad (cada
requisito tiene su mecanismo de acceso documentado en la matriz) ni la
cobertura de pruebas. Pero el dato más honesto para RQ4 no es un número
de cobertura -- es que la estrategia híbrida generó \textbf{fricción real
de implementación}: el equipo tuvo que abandonar
\texttt{@Procedure}/\texttt{@NamedStoredProcedureQuery} (el mecanismo
JPA formalmente diseñado para invocar procedimientos almacenados) y
migrar a \texttt{@Query} nativa con parámetros bindeados, por un defecto
documentado de la combinación Hibernate~6.2+/7.x con procedimientos
multi-\texttt{OUT} en PostgreSQL
(\texttt{docs/mediciones/backend/2026-07-28-fallo-invocacion-sp-multi-out.md},
\texttt{spring-projects/spring-data-jpa\#3393}, ya detallado en el
\autoref{cap:implementacion}). Es decir: la arquitectura híbrida no fue
gratis -- costó un ciclo de depuración real y un cambio de mecanismo de
invocación no planeado originalmente -- y ese costo es tan parte de la
respuesta a RQ4 como el resultado final positivo de cobertura y
trazabilidad.

\section{Comparación con trabajos relacionados}
\label{sec:disc-trabajos-relacionados}

El \autoref{cap:trabajos-relacionados} identificó una brecha concreta en
la literatura revisada: ningún trabajo de los 10 sistemas primarios
comparados combina, en un solo sistema, autenticación JWT con roles,
catálogo con autocompletado de ISBN, un chatbot con IA generativa real
usado tanto para recomendaciones como para consultas, y una arquitectura
híbrida de acceso a datos con evidencia empírica de rendimiento
reproducible. Con los resultados reales de este documento en mano, esa
brecha se sostiene y se refuerza con evidencia, no solo con ausencia de
contraejemplos: \citet{colley2018ormperformance} identifica el problema
de rendimiento que motiva una arquitectura híbrida ORM/SP pero no llega a
implementarla ni medirla; SGB-SaaS sí la implementa y la mide (RQ4,
arriba) -- incluyendo el costo real de implementarla, un dato que ese
trabajo tampoco podría reportar por no haber construido el sistema.
\citet{kusuma2017rediscaching} mide el efecto de Redis con métricas de
uso de CPU y tráfico, pero sin la evaluación estadística formal (Wilcoxon
pareado, Cliff's delta) que este documento aporta para su propio caso de
uso de caché (RQ1, arriba) -- aunque, con honestidad, esa misma
evaluación formal es la que expuso la limitación metodológica de
\texttt{cache\_frio} que un enfoque menos riguroso no habría descubierto.
Ningún trabajo comparado en el \autoref{cap:trabajos-relacionados} reporta una auditoría de
seguridad control por control con hallazgos reales corregidos
verificablemente (RQ2) ni una matriz de trazabilidad de 43 requisitos con
patrón de acceso a datos declarado por requisito (RQ4) -- estas dos
piezas de evidencia son, hasta donde este acervo bibliográfico permite
afirmar, una contribución metodológica de este proyecto más allá del
propio sistema.

\section{Hallazgos inesperados}
\label{sec:disc-hallazgos}

Un indicador indirecto pero concreto de que la verificación de este
proyecto fue real -- no un ejercicio cosmético para llenar una guía -- es
la cantidad de defectos genuinamente inesperados que el propio proceso de
desarrollo y auditoría encontró y corrigió, ninguno de ellos sembrado a
propósito para ``encontrar algo que reportar''. La comprobación del DSL
de Structurizr detectó que \texttt{structurizr/cli} es, en el estado
actual de ese ecosistema, una imagen retirada que imprime un aviso de
desuso pero termina con código de salida~$0$ para cualquier entrada
--válida o no--, por lo que una validación histórica contra esa imagen
nunca habría sido una validación real
(\autoref{cap:diseno-arquitectura}). Un defecto en la generación de
\texttt{credencial\_qr\_token} obligó a que Postgres generara el valor
vía \texttt{DEFAULT} de columna en vez de que el backend lo insertara
explícitamente (commit \texttt{1f1008e}). Un \texttt{MailHealthIndicator}
colgaba el \emph{healthcheck} del contenedor cuando no había SMTP
configurado, un defecto de configuración que solo se manifestaba en
entornos sin correo real (commit \texttt{14992d6}). Una
\texttt{NoResourceFoundException} devolvía \texttt{500} genérico en vez
de \texttt{404} -- el mismo hallazgo que cerró la verificación de A05 en
el \autoref{cap:resultados} (commit \texttt{951fae5}). Y una migración de
Flyway (\texttt{V3\_\_multas\_multiples\_por\_prestamo.sql}, entre otras)
fallaba en un despliegue limpio (\texttt{docker compose down -v} +
reconstrucción completa) porque la fuente de migraciones estaba duplicada
entre \texttt{database/migrations/} y \texttt{src/main/resources/db/migration/},
y \texttt{baseline-version} no reflejaba la versión más alta ya aplicada
al esquema (commit \texttt{03821d7}). Ninguno de estos 5 hallazgos estaba
en el plan de trabajo original -- todos surgieron de ejecutar el sistema
de verdad, contra Docker real, en vez de razonar sobre el código sin
correrlo. Esa es, en sí misma, la evidencia más directa de que el
programa de calidad de este proyecto (\autoref{cap:materiales-metodos})
no fue cosmético.

\section{Implicaciones prácticas}
\label{sec:disc-implicaciones}

Para una biblioteca universitaria real que evaluara adoptar un sistema
como SGB-SaaS, la lectura honesta de este capítulo es la siguiente: el
sistema cumple sus umbrales de rendimiento y cobertura con margen real, y
su superficie de seguridad auditada (6 de los controles OWASP más
relevantes para una aplicación con autenticación y datos personales) no
llegó limpia a la primera revisión -- tuvo defectos reales que se
corrigieron, lo cual es una señal de proceso saludable, no una alarma, si
se compara contra la alternativa de no auditar en absoluto. Lo que
todavía no está listo para una decisión de adopción informada es la
evidencia de usabilidad percibida por usuarios reales (RQ3, pendiente) y
la verificación de TLS end-to-end en un entorno de despliegue público
real (A02, parcial) -- ambas son, precisamente, las piezas que dependen
de infraestructura o de participantes externos al propio código, no de
trabajo de ingeniería adicional dentro del repositorio. Una institución
que considere este sistema debería tratar la Entrega Final como una base
arquitectónica y de calidad de código sólida, pero condicionar cualquier
decisión de producción a que esas dos piezas -- SUS y TLS real -- se
completen primero.
